% !TeX program = pdflatex
\documentclass[11pt,a4paper]{article}
\usepackage{../../recursos/latex/estilo-notas}

\begin{document}
\cabecalhoaula{02}{Regressão Linear e Regularização}{AME cap.~2--3; ISLP cap.~3 e~6}

%==============================================================================
\section*{Objetivos da aula}
%==============================================================================
Ao final desta aula o aluno deve ser capaz de:
\begin{itemize}[leftmargin=1.4cm]
  \item escrever o modelo de regressão linear em forma matricial e deduzir o
        estimador de mínimos quadrados (MQO);
  \item explicar por que MQO falha (ou nem existe) em \emph{altas dimensões}
        ($p$ grande, ou $p>n$);
  \item descrever seleção de subconjuntos e regressão \emph{stepwise};
  \item definir as penalizações \textbf{Ridge} ($\ell_2$) e \textbf{Lasso}
        ($\ell_1$) e interpretá-las como controle do balanço viés--variância;
  \item explicar \emph{por que o Lasso zera coeficientes} (esparsidade) e o Ridge
        não, com a intuição geométrica;
  \item reconhecer a interpretação bayesiana das duas penalizações.
\end{itemize}

\begin{emsala}
Conexão com a Aula~01: regressão linear é o método \emph{paramétrico} mais simples
--- baixa variância, possivelmente alto viés. Regularização é a ferramenta para
\emph{mover o ponto} no balanço viés--variância sem trocar de família de modelos.
Pergunta de abertura: \emph{``se eu tenho 50 covariáveis e 30 observações, o que
acontece com os mínimos quadrados?''}
\end{emsala}

%==============================================================================
\section{O modelo de regressão linear}
%==============================================================================
Métodos \textbf{paramétricos} supõem que a função de predição é descrita por um
número finito de parâmetros. A regressão linear usa
\begin{equation}\label{eq:linear}
  g(\x) = \bbeta^\top \x = \beta_0 x_0 + \beta_1 x_1 + \dots + \beta_p x_p,
  \qquad x_0 \equiv 1,
\end{equation}
com $\bbeta=(\beta_0,\dots,\beta_p)$. É essencial notar que $x_j$ \emph{não} precisa
ser uma covariável original: podemos criar atributos derivados ($x_j^2$, $x_jx_k$,
$\log x_j$, etc.). Assim, ``linear'' refere-se à linearidade \textbf{nos
parâmetros} $\bbeta$, não nas variáveis --- o modelo é muito mais flexível do que
parece à primeira vista.

\subsection{Mínimos quadrados ordinários (MQO)}
Estimamos $\bbeta$ minimizando a soma dos quadrados dos resíduos (RSS):
\begin{equation}\label{eq:mqo}
  \widehat{\bbeta}
   = \argmin_{\bbeta} \sum_{i=1}^n
     \big(Y_i - \bbeta^\top \X_i\big)^2
   = \argmin_{\bbeta} \norm{\bm{Y} - \bm{X}\bbeta}^2 ,
\end{equation}
onde $\bm{X}\in\R^{n\times(p+1)}$ é a matriz de delineamento (linha $i$ igual a
$\X_i^\top$, com a coluna de $1$s) e $\bm{Y}=(Y_1,\dots,Y_n)^\top$.

\begin{proposicao}[Equações normais]\label{prop:normal}
Se $\bm{X}^\top\bm{X}$ é invertível, o minimizador de \eqref{eq:mqo} é
\begin{equation}\label{eq:betahat}
  \widehat{\bbeta} = (\bm{X}^\top\bm{X})^{-1}\bm{X}^\top \bm{Y}.
\end{equation}
\end{proposicao}
\begin{proof}
Seja $S(\bbeta)=\norm{\bm{Y}-\bm{X}\bbeta}^2
=\bm{Y}^\top\bm{Y}-2\bbeta^\top\bm{X}^\top\bm{Y}+\bbeta^\top\bm{X}^\top\bm{X}\bbeta$.
O gradiente é $\nabla S(\bbeta)=-2\bm{X}^\top\bm{Y}+2\bm{X}^\top\bm{X}\bbeta$.
Igualando a zero obtêm-se as \emph{equações normais}
$\bm{X}^\top\bm{X}\,\bbeta=\bm{X}^\top\bm{Y}$; como a Hessiana
$2\bm{X}^\top\bm{X}$ é semidefinida positiva, o ponto crítico é mínimo. Se
$\bm{X}^\top\bm{X}$ é invertível, vale \eqref{eq:betahat}.
\end{proof}

\begin{observacao}
$\widehat{\bbeta}$ é o estimador de máxima verossimilhança sob o modelo
$Y_i = \bbeta^\top\X_i + \varepsilon_i$, com $\varepsilon_i\sim N(0,\sigma^2)$
i.i.d. A predição é $\rhat(\x)=\widehat{\bbeta}^\top\x$. Sob o modelo linear
correto, MQO é o melhor estimador linear não-viesado (Gauss--Markov).
\end{observacao}

\begin{observacao}[Quando a linearidade falha]
Mesmo quando $r(\x)$ não é linear, MQO estima a \emph{melhor aproximação linear} de
$r$ (a projeção de $Y$ no espaço gerado pelas covariáveis). Por isso, ampliar o
conjunto de atributos (§4.1, bases) é uma forma de reduzir esse viés de
especificação.
\end{observacao}

%==============================================================================
\section{O problema das altas dimensões}
%==============================================================================
Quando $p$ é grande relativamente a $n$, o MQO sofre de três males:
\begin{enumerate}[leftmargin=1.6cm]
  \item \textbf{Variância alta.} Quanto mais parâmetros, maior $\Var(\widehat{\bbeta})$
        --- o modelo superajusta (Aula~01). A matriz $\Var(\widehat{\bbeta})
        =\sigma^2(\bm{X}^\top\bm{X})^{-1}$ ``explode'' quando as colunas de
        $\bm{X}$ são quase colineares.
  \item \textbf{Não unicidade.} Se $p>n$, então $\bm{X}^\top\bm{X}$ é singular e
        \eqref{eq:betahat} \emph{não tem solução única}: há infinitos $\bbeta$ com
        RSS nulo (superajuste perfeito).
  \item \textbf{Interpretabilidade.} Muitas covariáveis irrelevantes dificultam a
        leitura do modelo.
\end{enumerate}
A resposta é \textbf{reduzir a complexidade efetiva}: ou selecionando um
subconjunto de covariáveis, ou \emph{encolhendo} os coeficientes (regularização).

\subsection{Seleção de subconjuntos e \emph{stepwise}}
\begin{description}[leftmargin=2.4cm,style=nextline]
  \item[Melhor subconjunto] ajusta-se o modelo para \emph{todos} os $2^p$
        subconjuntos de covariáveis e escolhe-se o melhor segundo um critério que
        penaliza complexidade (AIC, BIC) ou via validação cruzada. Inviável para
        $p$ moderado ($2^p$ explode).
  \item[\emph{Forward stepwise}] começa com o modelo vazio e adiciona, a cada
        passo, a covariável que mais reduz o RSS, até um critério de parada.
  \item[\emph{Backward stepwise}] começa com o modelo completo e remove, a cada
        passo, a covariável menos útil (exige $n>p$).
\end{description}

\begin{atencao}
Critérios como AIC/BIC ($\riscohat \approx \mathrm{EQM} + P(g)$, com
$P(g)\propto p/\widehat\sigma^2$) estimam o risco \emph{corrigindo} o otimismo do
erro de treino. Mas seleção \emph{stepwise} é gananciosa (não examina todos os
modelos) e a inferência clássica ``pós-seleção'' fica inválida: os $p$-valores
reportados ignoram que o modelo foi escolhido olhando os dados.
\end{atencao}

%==============================================================================
\section{Regularização: encolhendo coeficientes}
%==============================================================================
A ideia da regularização (ou \emph{penalização}, ou \emph{shrinkage}) é adicionar
ao RSS um termo que penaliza coeficientes grandes, desencorajando o superajuste:
\[
  \widehat{\bbeta}_{\text{pen}}
  = \argmin_{\bbeta}\Big\{ \norm{\bm{Y}-\bm{X}\bbeta}^2
       + \lambda\, \mathrm{Pen}(\bbeta) \Big\}.
\]
O hiperparâmetro $\lambda\ge 0$ controla a intensidade: $\lambda=0$ recupera MQO;
$\lambda\to\infty$ força $\bbeta\to 0$. \textbf{$\lambda$ é escolhido por validação
cruzada} (Aula~03).

\begin{ideia}
Regularizar = introduzir \emph{viés} de propósito para reduzir \emph{variância}.
Sob $p$ grande, o ganho na variância supera a perda no viés e o risco total cai.
É a Aula~01 em ação.
\end{ideia}

\subsection{Regressão Ridge ($\ell_2$)}
\begin{equation}\label{eq:ridge}
  \widehat{\bbeta}^{\,\text{ridge}}
   = \argmin_{\bbeta}\Big\{ \norm{\bm{Y}-\bm{X}\bbeta}^2
       + \lambda \sum_{j=1}^{p}\beta_j^2 \Big\}.
\end{equation}
(O intercepto $\beta_0$ em geral não é penalizado.) A grande vantagem é a solução
fechada:
\begin{equation}\label{eq:ridgesol}
  \widehat{\bbeta}^{\,\text{ridge}}
   = (\bm{X}^\top\bm{X} + \lambda \bm{I})^{-1}\bm{X}^\top\bm{Y}.
\end{equation}

\begin{proposicao}[Ridge sempre tem solução única]
Para todo $\lambda>0$, a matriz $\bm{X}^\top\bm{X}+\lambda\bm{I}$ é definida
positiva (logo invertível), mesmo quando $p>n$.
\end{proposicao}
\begin{proof}
Para qualquer $\bm{v}\neq 0$,
$\bm{v}^\top(\bm{X}^\top\bm{X}+\lambda\bm{I})\bm{v}
=\norm{\bm{X}\bm{v}}^2+\lambda\norm{\bm{v}}^2\ge \lambda\norm{\bm{v}}^2>0$.
\end{proof}
Eis a primeira virtude do Ridge: \emph{estabiliza} o problema mal-condicionado das
altas dimensões. Os coeficientes são encolhidos suavemente em direção a zero, mas
em geral \textbf{nenhum é exatamente zero}.

\subsection{O Lasso ($\ell_1$)}
\begin{equation}\label{eq:lasso}
  \widehat{\bbeta}^{\,\text{lasso}}
   = \argmin_{\bbeta}\Big\{ \norm{\bm{Y}-\bm{X}\bbeta}^2
       + \lambda \sum_{j=1}^{p}\abs{\beta_j} \Big\}.
\end{equation}
Trocamos a norma $\ell_2$ pela $\ell_1$. Não há solução fechada (o termo
$\abs{\beta_j}$ não é diferenciável em $0$), mas o problema é convexo e resolvido
eficientemente (\emph{coordinate descent}, LARS). A propriedade marcante:

\begin{ideia}
O Lasso faz \textbf{seleção de variáveis}: para $\lambda$ grande o suficiente,
vários $\widehat\beta_j$ são \emph{exatamente} zero. Obtemos um modelo
\emph{esparso} e interpretável de graça --- ao contrário do Ridge.
\end{ideia}

\subsubsection*{Por que o Lasso zera coeficientes e o Ridge não?}
Ambos têm a formulação equivalente \emph{com restrição}:
\[
  \min_{\bbeta}\ \norm{\bm{Y}-\bm{X}\bbeta}^2
  \quad\text{sujeito a}\quad
  \underbrace{\textstyle\sum_j\abs{\beta_j}\le t}_{\text{Lasso}}
  \quad\text{ou}\quad
  \underbrace{\textstyle\sum_j\beta_j^2\le t}_{\text{Ridge}}.
\]
Geometricamente, minimizamos uma forma quadrática (elipses de RSS centradas em
$\widehat{\bbeta}^{\text{MQO}}$) restrita a uma região. A região do Ridge é uma
\emph{bola} (suave: a solução quase nunca toca um eixo); a do Lasso é um
\emph{losango} com \emph{quinas sobre os eixos} --- e as elipses tendem a tocar a
restrição exatamente numa quina, onde algum $\beta_j=0$. Daí a esparsidade.

\begin{exemplo}[Caso ortonormal: fórmulas explícitas]
Quando $\bm{X}^\top\bm{X}=\bm{I}$, os estimadores têm forma fechada em função do
MQO $\widehat\beta_j$:
\[
  \widehat\beta_j^{\,\text{ridge}} = \frac{\widehat\beta_j}{1+\lambda},
  \qquad
  \widehat\beta_j^{\,\text{lasso}}
     = \operatorname{sinal}(\widehat\beta_j)\,\big(\abs{\widehat\beta_j}-\tfrac{\lambda}{2}\big)_+ .
\]
O Ridge \emph{multiplica} por um fator $<1$ (encolhe proporcionalmente, nunca
zera); o Lasso \emph{subtrai} uma constante e trunca em zero (\emph{soft
thresholding}) --- por isso anula os pequenos.
\end{exemplo}

\subsection{Interpretação bayesiana}
Ambas as penalizações são o \emph{logaritmo de uma priori} sobre $\bbeta$:
\begin{itemize}[leftmargin=1.4cm]
  \item Ridge $\equiv$ MAP com priori \textbf{gaussiana}
        $\beta_j\sim N(0,\tau^2)$: $-\log p(\bbeta)\propto \sum_j\beta_j^2$.
  \item Lasso $\equiv$ MAP com priori \textbf{de Laplace} (dupla exponencial)
        $p(\beta_j)\propto e^{-\abs{\beta_j}/b}$: $-\log p(\bbeta)\propto
        \sum_j\abs{\beta_j}$.
\end{itemize}
A priori de Laplace é mais ``pontuda'' em zero e tem caudas mais pesadas: ela
favorece soluções com muitos coeficientes \emph{nulos} e alguns grandes --- exatamente
a esparsidade observada.

\begin{atencao}
Como as penalizações somam $\beta_j^2$ ou $\abs{\beta_j}$, elas \textbf{dependem da
escala} das covariáveis. \emph{Sempre padronize} (média 0, variância 1) antes de
aplicar Ridge/Lasso, caso contrário variáveis em unidades grandes são penalizadas
de menos. Isso conecta diretamente com a Aula~07 (\emph{pipelines} e
\texttt{StandardScaler}).
\end{atencao}

%==============================================================================
\section{Prática em Python}
%==============================================================================
Em \texttt{scikit-learn}, as três variantes seguem a mesma interface; o
\texttt{alpha} é o nosso $\lambda$. Note a padronização \emph{dentro} do
\emph{pipeline} (Aula~07) para evitar vazamento.
\begin{lstlisting}
from sklearn.linear_model import LinearRegression, Ridge, Lasso, LassoCV
from sklearn.preprocessing import StandardScaler
from sklearn.pipeline import make_pipeline

# Lasso com lambda escolhido por validacao cruzada
modelo = make_pipeline(StandardScaler(),
                       LassoCV(cv=5, random_state=0))
modelo.fit(X_tr, y_tr)

lasso = modelo.named_steps["lassocv"]
print("lambda escolhido:", lasso.alpha_)
print("coeficientes nao-nulos:", (lasso.coef_ != 0).sum(), "de", X.shape[1])
\end{lstlisting}
O exemplo clássico (dados de câncer de próstata, ou as resenhas da Amazon no [AME])
mostra o padrão típico: o Lasso seleciona poucas covariáveis com desempenho
comparável ou superior ao MQO, especialmente quando $p$ é grande.

O notebook \texttt{Aula prática 02} percorre tudo isso no
\texttt{superconductivity.csv} ($p=81$), inclusive a pegadinha de escala entre
\texttt{Ridge} e \texttt{Lasso}.

%==============================================================================
\section*{Resumo}
%==============================================================================
\begin{itemize}[leftmargin=1.4cm]
  \item Regressão linear: $g(\x)=\bbeta^\top\x$, linear \emph{nos parâmetros};
        MQO $\widehat{\bbeta}=(\bm{X}^\top\bm{X})^{-1}\bm{X}^\top\bm{Y}$.
  \item Em altas dimensões MQO tem variância alta e nem existe se $p>n$.
  \item \textbf{Ridge} ($\ell_2$): solução fechada, sempre única, encolhe sem
        zerar.
  \item \textbf{Lasso} ($\ell_1$): faz seleção de variáveis (esparsidade) pela
        geometria das quinas / \emph{soft thresholding}.
  \item Ridge $\equiv$ priori gaussiana; Lasso $\equiv$ priori de Laplace.
  \item $\lambda$ é \emph{tuning parameter} (validação cruzada, Aula~03); sempre
        \emph{padronize} as covariáveis.
\end{itemize}

\paragraph{Leitura recomendada.} [AME] Capítulos~2 e~3 (§3.3 Lasso, §3.4 Ridge,
§3.6 interpretação bayesiana); [ISLP] Capítulo~3 (regressão linear) e Capítulo~6
(§6.1 seleção de subconjuntos, §6.2 \emph{shrinkage}).

\end{document}
